Una vez realizada la partición del conjunto de datos en subconjunto de entrenamiento y subconjunto de evaluación, se procedió al entrenamiento del clasificador. De los modelos descritos en el capítulo \ref{CAP:EST_ARTE}, se escogió un modelo \texttt{Random Forest}, debido a su capacidad de manejar conjuntos de datos grandes y complejos, su robustez frente a la variabilidad de los datos y su facilidad de uso. Además, como la primera parte de este trabajo fue definir nuevas características para el conjunto de datos, este modelo permite obtener la importancia de las características definidas en el conjunto de datos a la hora de clasificar el tráfico contenido en el subconjunto de evaluación. Los hiperparámetros definidos para el entrenamiento son los siguientes:

\PythonCode{Inicialización de Random Forest}
{Inicialización del clasificador \textit{Random Forest} con los hiperparámetros utilizados en los experimentos.}
{python/hiperparametros.py}{1}{20}{1}

El modelo se configuró con 400 árboles de decisión (\texttt{n_estimators = 400}). Cuantos más árboles se creen, más estable será la estimación, aunque tendrá un mayor coste computacional. En este caso se selecciono un valor medianamente elevado para que el clasificador fuera lo suficientemente robusto sin tener un coste muy elevado. En cada división de los árboles se empleó \texttt{max_features = "sqrt"}, lo cual significa que en cada división se consideran un total de $\sqrt{\text{Num\_Características}}$ aleatorias, aumentando la variabilidad entre los árboles y reduciendo la correlación entre estos. 

Para controlar la complejidad de los árboles se definieron \texttt{min_samples_split = 10}, que establece el mínimo de muestras para dividir un nodo y \texttt{min_samples_leaf = 2}, que establece el número mínimo de muestras que debe contener una hoja del árbol. De forma complementaria, se establece que \texttt{max_depth = None}, es decir, los árboles no estan limitados a una profundiad máxima, por lo que serán las otras dos reestricciones las que limiten el crecimiento de cada árbol.
